\chapter{Implementation and Technical Details}
\label{ch:implementation}

\section{Development Environment}

\subsection{Hardware Configuration}

The implementation was developed and evaluated on the following hardware:

\begin{itemize}
    \item \textbf{GPU}: Apple Silicon Metal Performance Shaders (MPS) / CUDA-compatible GPU
    \item \textbf{Memory}: 16 GB+ system RAM
    \item \textbf{Storage}: SSD with 50GB+ free space for dataset and checkpoints
    \item \textbf{Compute}: Multi-core CPU for data loading and preprocessing
\end{itemize}

The implementation selects the most efficient available accelerator in preference order: MPS $\rightarrow$ CUDA $\rightarrow$ CPU.

\subsection{Software Stack}

Key dependencies and their versions:

\begin{table}[H]
\centering
\caption{Software Dependencies}
\label{tab:dependencies}
\begin{tabular}{|l|c|l|}
\hline
\textbf{Library} & \textbf{Version} & \textbf{Purpose} \\
\hline
Python & 3.8+ & Programming language \\
PyTorch & 1.12+ & Deep learning framework \\
torchvision & 0.13+ & Image processing utilities \\
timm & 0.6+ & Pre-trained model library (Swin implementations) \\
scikit-learn & 1.0+ & Machine learning metrics \\
NumPy & 1.21+ & Numerical computations \\
Pandas & 1.3+ & Data manipulation \\
Matplotlib & 3.4+ & Visualization \\
PIL/Pillow & 8.0+ & Image I/O \\
tqdm & 4.60+ & Progress bars \\
\hline
\end{tabular}
\end{table}

\section{Model Implementation}

\subsection{Model Loading from Pre-trained Checkpoint}

The Swin Transformer is loaded from the \texttt{timm} library with ImageNet pre-training:

\begin{lstlisting}[language=Python, caption=Loading Swin Transformer Model]
import timm
import torch.nn as nn

MODEL_NAME = "swin_tiny_patch4_window7_224"

# Load pre-trained model
model = timm.create_model(
    MODEL_NAME, 
    pretrained=True, 
    num_classes=7
)

# Move to device
model = model.to(DEVICE)
\end{lstlisting}

The \texttt{num\_classes} parameter automatically adapts the classification head for 7 skin lesion categories.

\subsection{Custom Model Wrapper}

A wrapper class encapsulates the model with additional functionality for evaluation and uncertainty quantification:

\begin{lstlisting}[language=Python, caption=Custom Model Wrapper]
class SkinLesionClassifier(nn.Module):
    def __init__(self, model_name, num_classes=7):
        super().__init__()
        self.model = timm.create_model(
            model_name, 
            pretrained=True, 
            num_classes=num_classes
        )
        
    def forward(self, x):
        logits = self.model(x)
        probs = torch.softmax(logits, dim=1)
        return logits, probs
    
    def predict(self, x):
        _, probs = self.forward(x)
        predictions = torch.argmax(probs, dim=1)
        confidences = torch.max(probs, dim=1)[0]
        return predictions, confidences
\end{lstlisting}

\section{Data Loading Pipeline}

\subsection{Custom Dataset Class}

A PyTorch Dataset subclass handles image loading and augmentation:

\begin{lstlisting}[language=Python, caption=Custom Dataset Implementation]
from torch.utils.data import Dataset
from torchvision import transforms
from PIL import Image

class SkinLesionDataset(Dataset):
    def __init__(self, image_paths, labels, transform=None):
        self.image_paths = image_paths
        self.labels = labels
        self.transform = transform
    
    def __len__(self):
        return len(self.image_paths)
    
    def __getitem__(self, idx):
        image = Image.open(self.image_paths[idx]).convert('RGB')
        
        if self.transform:
            image = self.transform(image)
        
        label = self.labels[idx]
        return image, label
\end{lstlisting}

\subsection{Data Augmentation Pipeline}

Training data augmentation is implemented using torchvision transforms:

\begin{lstlisting}[language=Python, caption=Training Augmentation Pipeline]
from torchvision import transforms

train_transform = transforms.Compose([
    transforms.RandomRotation(180),
    transforms.RandomHorizontalFlip(p=0.5),
    transforms.RandomVerticalFlip(p=0.5),
    transforms.ColorJitter(
        brightness=0.2,
        contrast=0.2,
        saturation=0.2,
        hue=0.1
    ),
    transforms.RandomAffine(
        degrees=0,
        translate=(0.1, 0.1),
        scale=(0.8, 1.2)
    ),
    transforms.GaussianBlur(kernel_size=3),
    transforms.Resize((224, 224)),
    transforms.ToTensor(),
    transforms.Normalize(
        mean=[0.485, 0.456, 0.406],
        std=[0.229, 0.224, 0.225]
    )
])

val_transform = transforms.Compose([
    transforms.Resize((224, 224)),
    transforms.ToTensor(),
    transforms.Normalize(
        mean=[0.485, 0.456, 0.406],
        std=[0.229, 0.224, 0.225]
    )
])
\end{lstlisting}

\subsection{DataLoader Configuration}

Data is loaded using PyTorch's DataLoader with efficient batching:

\begin{lstlisting}[language=Python, caption=DataLoader Setup]
from torch.utils.data import DataLoader, WeightedRandomSampler

# Compute class weights
class_weights = compute_class_weights(train_labels)
sampler = WeightedRandomSampler(
    weights=class_weights, 
    num_samples=len(train_labels),
    replacement=True
)

# Create dataloaders
train_loader = DataLoader(
    train_dataset,
    batch_size=BATCH_SIZE,
    sampler=sampler,
    num_workers=NUM_WORKERS,
    pin_memory=True
)

val_loader = DataLoader(
    val_dataset,
    batch_size=BATCH_SIZE,
    shuffle=False,
    num_workers=NUM_WORKERS,
    pin_memory=True
)
\end{lstlisting}

\section{Training Loop Implementation}

\subsection{Training Procedure}

The main training loop implements standard practices for deep learning:

\begin{lstlisting}[language=Python, caption=Training Loop]
def train_epoch(model, train_loader, criterion, optimizer, device):
    model.train()
    total_loss = 0.0
    correct = 0
    total = 0
    
    progress_bar = tqdm(train_loader)
    for images, labels in progress_bar:
        images = images.to(device)
        labels = labels.to(device)
        
        # Forward pass
        optimizer.zero_grad()
        logits, probs = model(images)
        loss = criterion(logits, labels)
        
        # Backward pass
        loss.backward()
        torch.nn.utils.clip_grad_norm_(model.parameters(), max_norm=1.0)
        optimizer.step()
        
        # Metrics
        total_loss += loss.item()
        predictions = torch.argmax(probs, dim=1)
        correct += (predictions == labels).sum().item()
        total += labels.size(0)
        
        progress_bar.set_postfix({
            'loss': loss.item(),
            'acc': correct/total
        })
    
    return total_loss / len(train_loader), correct / total
\end{lstlisting}

\subsection{Validation Procedure}

Validation occurs without gradient computation:

\begin{lstlisting}[language=Python, caption=Validation Function]
def validate(model, val_loader, criterion, device):
    model.eval()
    total_loss = 0.0
    all_preds = []
    all_labels = []
    
    with torch.no_grad():
        for images, labels in val_loader:
            images = images.to(device)
            labels = labels.to(device)
            
            logits, probs = model(images)
            loss = criterion(logits, labels)
            
            total_loss += loss.item()
            predictions = torch.argmax(probs, dim=1)
            
            all_preds.extend(predictions.cpu().numpy())
            all_labels.extend(labels.cpu().numpy())
    
    accuracy = accuracy_score(all_labels, all_preds)
    return total_loss / len(val_loader), accuracy, all_preds, all_labels
\end{lstlisting}

\section{Loss Function Configuration}

\subsection{Weighted Cross-Entropy Loss}

The loss function weights classes inversely to their frequency:

\begin{lstlisting}[language=Python, caption=Loss Function Setup]
import torch.nn as nn

# Compute class weights
unique_classes, class_counts = np.unique(
    train_labels, 
    return_counts=True
)
class_weights = len(train_labels) / (len(unique_classes) * class_counts)
class_weights = torch.tensor(
    class_weights, 
    dtype=torch.float32
).to(DEVICE)

# Create loss function
criterion = nn.CrossEntropyLoss(weight=class_weights)
\end{lstlisting}

\section{Optimization Configuration}

\subsection{AdamW Optimizer with Learning Rate Scheduling}

\begin{lstlisting}[language=Python, caption=Optimizer and Scheduler Setup]
from torch.optim import AdamW
from torch.optim.lr_scheduler import CosineAnnealingLR

# Optimizer
optimizer = AdamW(
    model.parameters(),
    lr=LEARNING_RATE,
    weight_decay=WEIGHT_DECAY,
    betas=(0.9, 0.999)
)

# Learning rate scheduler
scheduler = CosineAnnealingLR(
    optimizer,
    T_max=NUM_EPOCHS,
    eta_min=1e-6
)
\end{lstlisting}

\section{Model Checkpointing and Early Stopping}

\subsection{Checkpoint Management}

The best model is saved based on validation accuracy:

\begin{lstlisting}[language=Python, caption=Checkpoint Management]
best_accuracy = 0.0
patience_counter = 0
patience = 5

for epoch in range(NUM_EPOCHS):
    train_loss, train_acc = train_epoch(...)
    val_loss, val_acc, _, _ = validate(...)
    
    print(f"Epoch {epoch+1}: Train Loss: {train_loss:.4f}, "
          f"Val Acc: {val_acc:.4f}")
    
    # Save best model
    if val_acc > best_accuracy:
        best_accuracy = val_acc
        torch.save({
            'epoch': epoch,
            'model_state_dict': model.state_dict(),
            'optimizer_state_dict': optimizer.state_dict(),
            'best_accuracy': best_accuracy
        }, CHECKPOINT_PATH)
        patience_counter = 0
    else:
        patience_counter += 1
    
    # Early stopping
    if patience_counter >= patience:
        print(f"Early stopping after {epoch+1} epochs")
        break
    
    scheduler.step()
\end{lstlisting}

\section{Evaluation Metrics Computation}

\subsection{Comprehensive Metrics}

Evaluation computes multiple metrics aligned with clinical requirements:

\begin{lstlisting}[language=Python, caption=Metrics Computation]
from sklearn.metrics import (
    accuracy_score, precision_score, recall_score, f1_score,
    roc_auc_score, confusion_matrix, classification_report
)

def compute_metrics(y_true, y_pred, y_proba):
    """Compute comprehensive evaluation metrics"""
    metrics = {
        'accuracy': accuracy_score(y_true, y_pred),
        'precision': precision_score(
            y_true, y_pred, 
            average='weighted', 
            zero_division=0
        ),
        'recall': recall_score(
            y_true, y_pred, 
            average='weighted', 
            zero_division=0
        ),
        'f1': f1_score(
            y_true, y_pred, 
            average='weighted', 
            zero_division=0
        ),
        'roc_auc': roc_auc_score(
            y_true, y_proba, 
            multi_class='ovr'
        )
    }
    return metrics
\end{lstlisting}

\section{Prediction and Uncertainty Quantification}

\subsection{Class Prediction with Confidence}

\begin{lstlisting}[language=Python, caption=Prediction with Confidence]
def predict_with_confidence(model, image_path, device):
    """Predict lesion class with confidence score"""
    image = Image.open(image_path).convert('RGB')
    image_tensor = val_transform(image).unsqueeze(0).to(device)
    
    with torch.no_grad():
        _, probs = model(image_tensor)
        confidence, prediction = torch.max(probs, dim=1)
    
    class_dict = {
        0: 'Actinic Keratosis',
        1: 'Basal Cell Carcinoma',
        2: 'Benign Keratosis-like Lesions',
        3: 'Dermatofibroma',
        4: 'Melanoma',
        5: 'Melanocytic Nevi',
        6: 'Vascular Lesion'
    }
    
    predicted_class = class_dict[prediction.item()]
    confidence_score = confidence.item()
    
    return predicted_class, confidence_score
\end{lstlisting}

\section{Reproducibility}

\subsection{Random Seed Management}

Reproducible results are ensured through consistent seed management:

\begin{lstlisting}[language=Python, caption=Seed Configuration]
def set_seed(seed=42):
    """Set random seeds for reproducibility"""
    random.seed(seed)
    np.random.seed(seed)
    torch.manual_seed(seed)
    torch.cuda.manual_seed_all(seed)
    torch.backends.cudnn.deterministic = True
    torch.backends.cudnn.benchmark = False

set_seed(42)
\end{lstlisting}

This implementation provides a complete, production-ready framework for training and evaluating the Swin Transformer-based skin lesion classifier.
